\chapter{Experimental Setup}
\label{ch:experiment}

This chapter describes how the configurations of Chapters~\ref{ch:framework} and
\ref{ch:ijtag} are turned into quantitative data. Section~\ref{sec:ex-flow}
describes the synthesis flow and tools. Section~\ref{sec:ex-nodes} describes the
three technology nodes. Section~\ref{sec:ex-matrix} defines the configuration
matrix. Section~\ref{sec:ex-metrics} defines the metrics and the normalisation
used throughout the results. Section~\ref{sec:ex-auto} describes the automation
that makes the sweep reproducible.

\section{Synthesis flow and tools}
\label{sec:ex-flow}

All configurations are functionally elaborated from the same VHDL-2008 source
base and synthesised with Synopsys Design Compiler~\cite{designcompiler}. For each configuration the
flow reads the technology library and the design packages, analyses and
elaborates the top-level accelerator with the chosen generics and redundancy
constants, compiles to the target standard-cell library, and writes a fixed set
of reports --- area, cell counts, references, resources, power, and timing.

Each configuration is synthesised against the same fixed, deliberately
aggressive clock constraint at a given node, so that the synthesis effort is
identical across all configurations. Because the target period is intentionally
tighter than the design can meet, every configuration reports a (negative) slack;
slack is therefore \emph{not} a useful discriminator here. Instead, the timing
figure of merit used throughout is the \emph{critical-path (data-arrival) delay}
of the longest path, which the tool reports regardless of slack. Since the same
constraint and effort are applied to every configuration at a node, the
data-arrival delays are directly comparable, and their ratio to the baseline
captures exactly how much each redundancy scheme lengthens the critical path. The
voters and comparators of DMR/TMR are thin combinational layers on the
functional-unit outputs, so they are expected to add little to the critical path
--- a hypothesis the results confirm.

\section{Technology nodes}
\label{sec:ex-nodes}

To test whether the cost of redundancy is stable across process generations, each
configuration is synthesised at three nodes spanning a wide range of maturity:

\begin{itemize}
  \item \textbf{IHP 130\,nm} (\texttt{sg13g2})~\cite{ihp} --- a mature, open
        process-design-kit node, representative of robust, cost-sensitive and
        radiation-tolerant designs.
  \item \textbf{NanGate 45\,nm}~\cite{nangate} --- an open, widely used academic
        standard-cell library representative of a mainstream node.
  \item \textbf{NanGate 15\,nm}~\cite{nangate} --- an open FinFET-era library representative of
        an advanced node.
\end{itemize}

These three points let the study distinguish effects that are intrinsic to the
\emph{architecture} of the redundancy schemes (and therefore stable across nodes)
from effects that depend on the \emph{technology}. A practical consequence of
mixing such different libraries is that their reporting units differ --- dynamic
power is reported in milliwatts at two nodes and microwatts at another, and the
advanced library reports timing and power on a different absolute scale. The
analysis handles this by reading the units from each report and, wherever
absolute values are compared across nodes, by \emph{normalising to the per-node
baseline} so that unit and scale differences cancel.

\section{Configuration matrix}
\label{sec:ex-matrix}

The full sweep is the Cartesian product of the configurations and the technology
nodes. The configurations comprise the two scenarios introduced in
Chapter~\ref{ch:framework}:

\begin{itemize}
  \item \textbf{Scenario~A (single core)} --- baseline; each of the four
        arithmetic units (adder, multiplier, shifter, comparator) hardened
        individually with DMR and with TMR; and the full datapath hardened with
        DMR and with TMR. This is eleven configurations.
  \item \textbf{Scenario~B (replicated)} --- replicated baseline; replicated with
        DMR voting; replicated with TMR voting. This is three configurations.
\end{itemize}

Fourteen configurations at three nodes give \textbf{42 synthesis runs} for the
redundancy study, complemented by the \textbf{18 instrumented runs} of
Chapter~\ref{ch:ijtag} (three sensor counts $\times$ three nodes $\times$ two
accounting modes). The redundancy matrix was already listed in
Table~\ref{tab:matrix}.

\section{Metrics and normalisation}
\label{sec:ex-metrics}

From each run the following metrics are extracted:

\begin{itemize}
  \item \textbf{Area.} Total cell area, decomposed into combinational,
        buffer/inverter, sequential (non-combinational), and net-interconnect
        area; plus the number of total, combinational, and sequential cells.
  \item \textbf{Power.} Switching, internal, and leakage power, and their total,
        all normalised to milliwatts by reading each report's unit declarations.
        Power is reported by Design Compiler's hierarchical power analysis with
        default switching activity, so the absolute values are estimates; the
        \emph{ratios} between configurations are the meaningful quantity.
  \item \textbf{Timing.} The critical-path (data-arrival) delay of the longest
        path, used in preference to slack as discussed in
        Section~\ref{sec:ex-flow}.
\end{itemize}

Throughout the results, the headline figure of merit is the \emph{overhead
factor} --- the metric of a hardened configuration divided by the metric of the
baseline \emph{at the same node}. Normalising in this way removes both the
absolute-scale differences between libraries and any common-mode contribution of
the surrounding core, isolating the cost attributable to redundancy itself.

\section{Automation and reproducibility}
\label{sec:ex-auto}

Running 42 synthesis jobs by hand, each requiring a precise combination of
generics and redundancy constants, would be both tedious and unreliable. The
sweep is therefore fully automated by a driver that, for each (node,
configuration) pair, edits the relevant source constants, invokes Design
Compiler, and files the resulting reports into a clearly named directory. The
two scenarios are handled by the same driver: for Scenario~B it toggles the
mapping-layer \texttt{DMR\_ENABLED}/\texttt{TMR\_ENABLED} constants with
replication enabled, and for Scenario~A it toggles the per-unit
\texttt{*\_DMR\_EN}/\texttt{*\_TMR\_EN} constants with replication disabled. The
source files it edits are backed up before any change and restored when the
sweep finishes --- even on error or interruption --- so the working tree is
always left clean. The data-reduction and plotting scripts described in
Appendix~\ref{app:tools} then parse every report into a single table and
generate the figures and tables of the next chapter directly from it, so that
the entire path from raw reports to thesis figures is reproducible by re-running
two scripts.
